\chapter{Selection and Panel Methods}
\label{ch:selection-panel}

\section{Introduction}

Sample selection occurs when outcomes are only observed for a non-random subsample. This chapter develops:
\begin{enumerate}
\item \textbf{Heckman selection correction}: Two-step estimator for selection bias
\item \textbf{Panel data methods}: Correlated random effects for unobserved heterogeneity
\end{enumerate}

Both methods address endogeneity from unobserved factors, but through different mechanisms.

\subsection{Motivating Example: Wage Equations}

Consider estimating returns to education:
\begin{itemize}
\item $Y$: Log wages (outcome)
\item $T$: Years of education (treatment)
\item $S$: Employment status (selection)
\end{itemize}

Problem: Wages are only observed for employed individuals. If unobserved factors (ability, motivation) affect both employment \textit{and} wages, OLS on the employed sample is biased.


\section{The Sample Selection Problem}

\subsection{Formal Setup}

Let $Y^*$ be the latent outcome and $S^*$ be the latent selection tendency:
\begin{align}
Y^* &= X'\beta + u \label{eq:outcome-latent}\\
S^* &= Z'\gamma + v \label{eq:selection-latent}
\end{align}

We observe:
\begin{itemize}
\item $S = \mathbf{1}\{S^* > 0\}$: Binary selection indicator
\item $Y = Y^* \cdot S$: Outcome only when selected
\end{itemize}

The key problem: If $\text{Cov}(u, v) \neq 0$, then:
\begin{equation}
E[Y \mid X, S=1] = X'\beta + E[u \mid v > -Z'\gamma] \neq X'\beta
\end{equation}

\subsection{OLS Inconsistency}

\begin{proposition}[Selection Bias]
Under selection with $\text{Corr}(u, v) = \rho \neq 0$:
\begin{equation}
\hat{\beta}_{\text{OLS}} \xrightarrow{p} \beta + \rho \sigma_u \cdot \frac{E[\lambda(Z'\gamma) \mid S=1, X]}{X}
\end{equation}
where $\lambda(\cdot)$ is the Inverse Mills Ratio.
\end{proposition}

\begin{insight}
Selection bias direction depends on $\rho$:
\begin{itemize}
\item $\rho > 0$: Selected individuals have higher $u$ (upward bias if $\lambda > 0$)
\item $\rho < 0$: Selected individuals have lower $u$ (downward bias)
\end{itemize}
For wage equations, $\rho > 0$ is common---employed individuals tend to have favorable unobserved characteristics.
\end{insight}


\section{Heckman Two-Step Correction}

\subsection{The Inverse Mills Ratio}

\textcite{heckman1979sample}'s insight: Under joint normality of $(u, v)$:
\begin{equation}
E[u \mid v > -Z'\gamma] = \rho \sigma_u \cdot \lambda(Z'\gamma)
\end{equation}

where the Inverse Mills Ratio is:
\begin{equation}
\lambda(a) = \frac{\phi(a)}{\Phi(a)}
\end{equation}
with $\phi(\cdot)$ and $\Phi(\cdot)$ the standard normal PDF and CDF.

\subsection{Two-Step Procedure}

\begin{algorithm}[Heckman Two-Step]
\label{alg:heckman}
\begin{enumerate}
\item \textbf{Step 1 (Selection equation)}: Estimate probit on full sample
\begin{equation}
P(S_i = 1 \mid Z_i) = \Phi(Z_i'\gamma)
\end{equation}
Obtain $\hat{\gamma}$ and predicted probabilities $\hat{p}_i = \Phi(Z_i'\hat{\gamma})$.

\item \textbf{Step 2 (Compute IMR)}: For selected observations, compute
\begin{equation}
\hat{\lambda}_i = \frac{\phi(\Phi^{-1}(\hat{p}_i))}{\hat{p}_i}
\end{equation}

\item \textbf{Step 3 (Outcome equation)}: On selected sample, estimate
\begin{equation}
Y_i = X_i'\beta + \lambda_{\text{coef}} \cdot \hat{\lambda}_i + \varepsilon_i
\end{equation}
The coefficient $\lambda_{\text{coef}} = \rho \sigma_u$ captures selection bias.

\item \textbf{Step 4 (Correct standard errors)}: Use Heckman-corrected variance:
\begin{equation}
\text{Var}(\hat{\beta}) = \sigma^2 (X'X)^{-1} + \rho^2 (X'X)^{-1}(X'\Delta X)(X'X)^{-1}
\end{equation}
where $\Delta$ depends on the selection equation uncertainty.
\end{enumerate}
\end{algorithm}

\subsection{Implementation}

\begin{minted}[fontsize=\small]{python}
from causal_inference.selection import heckman_two_step

result = heckman_two_step(
    outcome=wages,                    # NaN for unemployed
    selected=employed,                # 1 if employed, 0 otherwise
    selection_covariates=np.column_stack([education, age, married]),
    outcome_covariates=education.reshape(-1, 1),  # Exclusion: age, married
    alpha=0.05,
    add_intercept=True,
)

# Results
print(f"Education coefficient: {result['estimate']:.4f}")
print(f"SE (Heckman-corrected): {result['se']:.4f}")
print(f"95% CI: [{result['ci_lower']:.4f}, {result['ci_upper']:.4f}]")

# Selection parameters
print(f"\nSelection correlation (rho): {result['rho']:.4f}")
print(f"Lambda coefficient: {result['lambda_coef']:.4f}")
print(f"Lambda p-value: {result['lambda_pvalue']:.4f}")
\end{minted}


\section{Identification and Exclusion Restrictions}

\subsection{The Identification Problem}

Without additional structure, $\beta$ is not identified from functional form alone.

\begin{assumption}[Exclusion Restriction]
At least one variable in $Z$ (selection equation) is \textit{not} in $X$ (outcome equation).
\end{assumption}

\begin{itemize}
\item $Z$ includes variables affecting selection but not outcome
\item These \textit{excluded} variables provide identifying variation
\item Without exclusion, identification relies on IMR nonlinearity (fragile)
\end{itemize}

\subsection{Examples of Exclusion Restrictions}

\begin{table}[h]
\centering
\caption{Common Exclusion Restrictions}
\begin{tabular}{lll}
\toprule
Setting & Excluded Variable & Justification \\
\midrule
Wage equation & Spouse's income & Affects employment, not wages \\
Wage equation & Number of children & Affects labor supply, not wage offer \\
Migration & Distance to destination & Affects migration, not outcomes \\
Survey & Interview length & Affects response, not true answer \\
\bottomrule
\end{tabular}
\end{table}

\subsection{Diagnostics}

\begin{minted}[fontsize=\small]{python}
from causal_inference.selection import diagnose_identification

id_check = diagnose_identification(
    selection_covariates=Z,
    outcome_covariates=X,
    threshold=0.99,
)

print(f"Has exclusion restriction: {id_check['has_exclusion']}")
print(f"Excluded variable indices: {id_check['excluded_indices']}")
print(f"Identification strength: {id_check['identification_strength']}")
# 'strong' = good exclusion, 'weak' = high collinearity, 'fragile' = no exclusion
\end{minted}


\section{Testing for Selection Bias}

\subsection{The Lambda Test}

Test $H_0: \lambda_{\text{coef}} = 0$ (no selection bias):
\begin{equation}
t = \frac{\hat{\lambda}_{\text{coef}}}{\text{SE}(\hat{\lambda}_{\text{coef}})}
\end{equation}

\begin{minted}[fontsize=\small]{python}
from causal_inference.selection import selection_bias_test

test_result = selection_bias_test(
    lambda_coef=result['lambda_coef'],
    lambda_se=result['lambda_se'],
    alpha=0.05,
)

print(f"t-statistic: {test_result['statistic']:.3f}")
print(f"p-value: {test_result['pvalue']:.4f}")
print(f"Reject H0: {test_result['reject_null']}")
print(test_result['interpretation'])
\end{minted}

\subsection{Interpretation}

\begin{itemize}
\item \textbf{Reject $H_0$}: Significant selection bias; Heckman correction needed
\item \textbf{Fail to reject}: No evidence of selection bias; OLS may be valid
\end{itemize}

\begin{warning}
Failure to reject does \textit{not} prove absence of selection bias. The test may have low power, especially with:
\begin{itemize}
\item Small samples
\item Weak exclusion restrictions
\item Low selection rate
\end{itemize}
\end{warning}


\section{Panel Data Methods}

\subsection{Unobserved Heterogeneity}

Panel data provides repeated observations on the same units:
\begin{equation}
Y_{it} = X_{it}'\beta + \alpha_i + \varepsilon_{it}
\end{equation}
where $\alpha_i$ is the unit-specific unobserved effect.

The challenge: If $\text{Cov}(\alpha_i, X_{it}) \neq 0$, pooled OLS is inconsistent.

\subsection{Fixed Effects vs Random Effects}

\begin{table}[h]
\centering
\caption{Panel Estimator Comparison}
\begin{tabular}{lll}
\toprule
Approach & Assumption & Trade-off \\
\midrule
Fixed Effects (FE) & $\text{Cov}(\alpha_i, X_{it})$ unrestricted & Cannot estimate time-invariant effects \\
Random Effects (RE) & $\text{Cov}(\alpha_i, X_{it}) = 0$ & Efficient but inconsistent if violated \\
\bottomrule
\end{tabular}
\end{table}

\subsection{Correlated Random Effects (Mundlak)}

\textcite{mundlak1978pooling} proposes modeling the correlation:
\begin{equation}
\alpha_i = \bar{X}_i'\psi + c_i, \quad c_i \perp\!\!\!\perp X_{it}
\end{equation}

where $\bar{X}_i = \frac{1}{T_i}\sum_{t=1}^{T_i} X_{it}$ is the unit-specific mean.

\begin{proposition}[Mundlak Model]
The augmented model:
\begin{equation}
Y_{it} = X_{it}'\beta + \bar{X}_i'\psi + c_i + \varepsilon_{it}
\end{equation}
yields consistent $\hat{\beta}$ under FE assumptions while retaining RE efficiency.
\end{proposition}

\subsection{DML with Correlated Random Effects}

For causal inference with panel data, we combine CRE with double machine learning:

\begin{minted}[fontsize=\small]{python}
import numpy as np
from sklearn.model_selection import KFold
from sklearn.ensemble import RandomForestRegressor

def dml_cre_ate(Y, T, X, unit_ids, n_folds=5):
    """
    DML-CRE estimator for panel data ATE.

    Parameters
    ----------
    Y : array
        Outcomes (n_obs,)
    T : array
        Treatment indicator (n_obs,)
    X : array
        Time-varying covariates (n_obs, p)
    unit_ids : array
        Unit identifiers (n_obs,)

    Returns
    -------
    ate, se : float
        Estimated ATE and standard error
    """
    # Compute unit means for CRE
    unique_ids = np.unique(unit_ids)
    X_bar = np.zeros((len(Y), X.shape[1]))

    for uid in unique_ids:
        mask = unit_ids == uid
        X_bar[mask] = X[mask].mean(axis=0)

    # Augmented covariates: [X, X_bar]
    X_aug = np.column_stack([X, X_bar])

    # Cross-fitting
    kf = KFold(n_splits=n_folds, shuffle=True, random_state=42)
    resid_Y = np.zeros_like(Y)
    resid_T = np.zeros_like(T)

    for train_idx, test_idx in kf.split(Y):
        # Nuisance models
        m_Y = RandomForestRegressor(n_estimators=100, random_state=42)
        m_T = RandomForestRegressor(n_estimators=100, random_state=42)

        m_Y.fit(X_aug[train_idx], Y[train_idx])
        m_T.fit(X_aug[train_idx], T[train_idx])

        resid_Y[test_idx] = Y[test_idx] - m_Y.predict(X_aug[test_idx])
        resid_T[test_idx] = T[test_idx] - m_T.predict(X_aug[test_idx])

    # ATE via residual regression
    ate = np.sum(resid_Y * resid_T) / np.sum(resid_T ** 2)

    # Cluster-robust SE
    # (simplified - full implementation uses cluster bootstrap)
    residuals = resid_Y - ate * resid_T
    se = np.std(residuals) / np.sqrt(np.sum(resid_T ** 2))

    return ate, se
\end{minted}


\section{Combining Selection and Panel}

\subsection{Panel Selection Models}

When outcomes are observed only for selected observations in panel data:
\begin{align}
S_{it}^* &= Z_{it}'\gamma + \alpha_i^S + v_{it} \\
Y_{it}^* &= X_{it}'\beta + \alpha_i^Y + u_{it}
\end{align}

Challenges:
\begin{enumerate}
\item Selection equation has unit-specific effect $\alpha_i^S$
\item Correlation between $\alpha_i^S$ and $\alpha_i^Y$ creates additional selection channel
\item Standard Heckman assumes cross-sectional independence
\end{enumerate}

\subsection{Panel Heckman Estimator}

Extension to panel data:
\begin{enumerate}
\item Estimate random effects probit for selection
\item Compute panel-corrected IMR
\item Include unit means (CRE) in outcome equation
\item Cluster standard errors by unit
\end{enumerate}


\section{Implementation Details}

\subsection{Complete Heckman Workflow}

\begin{minted}[fontsize=\small]{python}
import numpy as np
from causal_inference.selection import (
    heckman_two_step,
    selection_bias_test,
    diagnose_identification,
)

# Simulate data with selection
np.random.seed(42)
n = 2000

# Covariates
education = np.random.randn(n) * 2 + 12
age = np.random.randn(n) * 10 + 35
married = np.random.binomial(1, 0.6, n)

# Latent selection (employment)
# Age and married affect employment but not wages (exclusion)
v = np.random.randn(n)
selection_latent = -2 + 0.1 * education + 0.05 * age + 0.3 * married + v
employed = (selection_latent > 0).astype(int)

# Latent wages
# Only education affects wages
u = np.random.randn(n)
# Correlation between u and v creates selection bias
u = 0.5 * v + np.sqrt(1 - 0.5**2) * np.random.randn(n)
wages = 10 + 0.5 * education + u

# Observed wages (only for employed)
wages_observed = np.where(employed == 1, wages, np.nan)

print(f"Employment rate: {employed.mean():.1%}")

# Check identification
Z = np.column_stack([education, age, married])
X = education.reshape(-1, 1)
id_check = diagnose_identification(Z, X)
print(f"Identification: {id_check['identification_strength']}")

# Estimate
result = heckman_two_step(
    outcome=wages_observed,
    selected=employed,
    selection_covariates=Z,
    outcome_covariates=X,
)

print(f"\n=== Heckman Results ===")
print(f"Education coef: {result['estimate']:.4f} (true: 0.50)")
print(f"SE: {result['se']:.4f}")
print(f"95% CI: [{result['ci_lower']:.4f}, {result['ci_upper']:.4f}]")
print(f"\nSelection rho: {result['rho']:.4f} (true: 0.50)")
print(f"Lambda coef: {result['lambda_coef']:.4f}")
print(f"Lambda p-value: {result['lambda_pvalue']:.4f}")

# Test for selection bias
test = selection_bias_test(result['lambda_coef'], result['lambda_se'])
print(f"\nSelection test: {test['interpretation']}")

# Compare to naive OLS
from sklearn.linear_model import LinearRegression
mask = employed == 1
ols = LinearRegression().fit(education[mask].reshape(-1,1), wages[mask])
print(f"\nNaive OLS: {ols.coef_[0]:.4f} (biased)")
\end{minted}


\section{Monte Carlo Validation}

\subsection{Test 1: Heckman Consistency}

\begin{table}[h]
\centering
\caption{Heckman Monte Carlo Results ($n=1000$, $\rho=0.5$, 1,000 runs)}
\begin{tabular}{lcccc}
\toprule
Estimator & True $\beta$ & Mean Est. & Bias & Coverage \\
\midrule
Naive OLS & 0.50 & 0.72 & 0.22 & 12.3\% \\
Heckman & 0.50 & 0.502 & 0.004 & 94.6\% \\
\bottomrule
\end{tabular}
\end{table}

\subsection{Test 2: Lambda Test Power}

\begin{table}[h]
\centering
\caption{Selection Test Power ($\alpha=0.05$, $n=500$)}
\begin{tabular}{lcc}
\toprule
True $\rho$ & Rejection Rate & Interpretation \\
\midrule
0.0 & 5.1\% & Correct size \\
0.2 & 34.2\% & Low power \\
0.5 & 89.7\% & Good power \\
0.8 & 99.8\% & Excellent power \\
\bottomrule
\end{tabular}
\end{table}


\section{Discussion}

\subsection{Key Takeaways}

\begin{enumerate}
\item \textbf{Selection bias}: Outcomes observed only for selected sample create bias
\item \textbf{Heckman correction}: Two-step procedure with Inverse Mills Ratio
\item \textbf{Exclusion restrictions}: Critical for identification
\item \textbf{Panel methods}: CRE addresses time-invariant unobserved heterogeneity
\item \textbf{Standard errors}: Must correct for two-stage estimation uncertainty
\end{enumerate}

\subsection{Practical Recommendations}

\begin{enumerate}
\item Always check for exclusion restrictions using diagnostics
\item Test for selection bias before assuming Heckman is needed
\item Report both OLS and Heckman estimates for comparison
\item For panel data, use CRE to retain time-invariant effects
\item Cluster standard errors appropriately (unit-level for panels)
\end{enumerate}

\subsection{Limitations}

\begin{itemize}
\item \textbf{Normality assumption}: Joint normality of $(u, v)$ may be violated
\item \textbf{Functional form}: Results sensitive to probit vs logit choice
\item \textbf{Exclusion validity}: Hard to verify exclusion restriction validity
\item \textbf{Generated regressor}: Two-step SEs require correction
\end{itemize}

\subsection{Connection to Other Methods}

\begin{itemize}
\item \textbf{IV}: Selection model is like IV with the IMR as instrument
\item \textbf{MTE} (Chapter~\ref{ch:mte}): MTE generalizes Heckman to allow heterogeneous selection
\item \textbf{Bounds}: When selection assumptions fail, partial identification provides bounds
\end{itemize}
